\chapter{Models}
\label{app:model-algorithms}

This appendix supplements the operating-principle summary in
Table~\ref{tab:model_roster} with model-specific pseudocode. The common
synthetic-gap and fixed-origin forecasting protocols remain defined in
Chapter~\ref{ChapModels}; they are not repeated here. Likewise,
Algorithm~\ref{alg:auxiliary-driver-gapfill} already documents auxiliary
environmental-driver reconstruction. The algorithms below therefore focus on
how each benchmark model learns or produces predictions. ``GF'' denotes
gap-filling and ``FC'' denotes forecasting.

\section{Multivariate Imputation}

MICE and HyperImpute treat FCH\textsubscript{4} as an incomplete column in a
joint feature--target matrix. MICE uses a fixed Bayesian-ridge conditional
model \cite{vanbuuren2011mice}, whereas HyperImpute selects a conditional
learner separately for each variable through internal validation
\cite{jarrett2022hyperimpute}.

\begin{algorithm}[H]
\caption{MICE reconstruction with Bayesian-ridge conditional models}
\label{alg:app-mice}
\small
\begin{algorithmic}[1]
\Require Pooled incomplete matrix $Z=[y,X,\text{tower indicators}]$;
held-out target positions $H$
\Ensure Reconstructed target $\hat y_H$
\State Set $Z_{H,y}$ to missing before constructing the imputer input
\State Initialise each incomplete column from its observed marginal values
\For{chained iteration $q\gets1$ to $K$}
    \For{each incomplete column $j$}
        \State Define observed rows $O_j$ and missing rows $M_j$
        \State Fit Bayesian ridge $f_j^{(q)}:Z_{O_j,-j}\mapsto Z_{O_j,j}$
        \State Update $Z_{M_j,j}\gets f_j^{(q)}(Z_{M_j,-j})$
    \EndFor
    \If{the iterative-imputer convergence rule is satisfied}
        \State Break
    \EndIf
\EndFor
\State Extract $\hat y_H\gets Z_{H,y}$ using the held-out rows' absolute positions
\State \Return $\hat y_H$
\end{algorithmic}
\end{algorithm}

\begin{algorithm}[H]
\caption{HyperImpute reconstruction with variable-specific AutoML models}
\label{alg:app-hyperimpute}
\small
\begin{algorithmic}[1]
\Require Pooled incomplete matrix $Z=[y,X,\text{tower indicators}]$;
held-out positions $H$; candidate learner library $\mathcal L$
\Ensure Reconstructed target $\hat y_H$
\State Set $Z_{H,y}$ to missing and initialise incomplete columns
\For{each chained-imputation round}
    \For{each incomplete column $j$}
        \State Determine whether $Z_j$ requires regression or classification
        \For{each compatible learner $f\in\mathcal L$}
            \State Estimate internal validation loss $L_j(f)$ on observed rows $O_j$
        \EndFor
        \State Select $f_j^*\gets\arg\min_{f\in\mathcal L}L_j(f)$
        \State Refit $f_j^*$ on $O_j$ and update
        $Z_{M_j,j}\gets f_j^*(Z_{M_j,-j})$
    \EndFor
    \If{the imputation sequence converges or reaches its iteration limit}
        \State Break
    \EndIf
\EndFor
\State \Return the completed target values $\hat y_H$
\end{algorithmic}
\end{algorithm}

\section{Tree-Based Models}

Random Forest, XGBoost and LightGBM were used in both tasks. For gap-filling,
the query rows are the deliberately masked hours. For forecasting, each model
is fitted at a forecast origin and its predicted FCH\textsubscript{4} value is
appended to the available target history before constructing the next day's
autoregressive features. Preprocessing is always estimated from permitted
training rows. The algorithms follow Breiman \cite{breiman2001randomforests},
Chen and Guestrin \cite{chen2016xgboost}, and Ke et al.\ \cite{ke2017lightgbm}.
The unweighted voting-regression ensemble then averages the three aligned
forecast chains following the forecast-combination principle of Bates and
Granger \cite{bates1969combination}.

\begin{algorithm}[H]
\caption{Random Forest reconstruction or recursive forecasting}
\label{alg:app-rf}
\small
\begin{algorithmic}[1]
\Require Training rows $(X_O,y_O)$; query rows $X_Q$; number of trees $B$;
minimum leaf size $m$; task $\tau\in\{\mathrm{GF},\mathrm{FC}\}$
\Ensure Predictions $\hat y_Q$
\State Fit mean imputation on $X_O$ and apply it to $X_O$ and $X_Q$
\State Use $(B,m)=(500,5)$ for GF and $(500,10)$ for FC
\For{$b\gets1$ to $B$}
    \State Draw a bootstrap sample of $(X_O,y_O)$
    \State Grow $T_b$ using squared-error reduction and random candidate features
    \State Stop a branch when further splitting violates the minimum leaf size
\EndFor
\State Define $f(x)\gets B^{-1}\sum_{b=1}^{B}T_b(x)$
\If{$\tau=\mathrm{GF}$}
    \State Predict the masked hourly rows directly: $\hat y_Q\gets f(X_Q)$
\Else
    \For{forecast day $h\gets1$ to $365$}
        \State Build day-$h$ features from known covariates and available target history
        \State Predict $\hat y_h\gets f(x_h)$ and append $\hat y_h$ to that history
    \EndFor
\EndIf
\State \Return $\hat y_Q$
\end{algorithmic}
\end{algorithm}

\begin{algorithm}[H]
\caption{XGBoost reconstruction or recursive forecasting}
\label{alg:app-xgboost}
\small
\begin{algorithmic}[1]
\Require Training rows $(X_O,y_O)$; query rows $X_Q$; boosting rounds $B$;
learning rate $\eta$
\Ensure Predictions $\hat y_Q$
\State Fit training-only predictor imputation and initialise $F_0(x)$ with a constant
\For{$b\gets1$ to $B$}
    \For{each $i\in O$}
        \State Compute $g_i=\partial_F\ell(y_i,F_{b-1}(x_i))$ and
        $h_i=\partial_F^2\ell(y_i,F_{b-1}(x_i))$
    \EndFor
    \State Grow $T_b$ by maximising regularised split gain from sums of $g_i$ and $h_i$
    \State Assign each leaf its regularised optimum weight
    \State Update $F_b(x)\gets F_{b-1}(x)+\eta T_b(x)$
\EndFor
\State Use 500 rounds and minimum child weight 5 for GF
\State Use 400 rounds, depth 2, $\eta=0.02$ and minimum child weight 10 for FC
\State Predict GF rows directly or generate the FC chain by appending each prediction
\State \Return $F_B(X_Q)$ or the resulting recursive chain
\end{algorithmic}
\end{algorithm}

\begin{algorithm}[H]
\caption{LightGBM reconstruction or recursive forecasting}
\label{alg:app-lightgbm}
\small
\begin{algorithmic}[1]
\Require Training rows $(X_O,y_O)$; query rows $X_Q$; boosting rounds $B$;
learning rate $\eta$
\Ensure Predictions $\hat y_Q$
\State Fit training-only predictor imputation and bin continuous predictors
\State Initialise $F_0(x)$ with a constant
\For{$b\gets1$ to $B$}
    \State Calculate loss gradients from $y_O$ and $F_{b-1}(X_O)$
    \State Accumulate gradient statistics by feature-histogram bin
    \While{a permissible positive-gain split exists}
        \State Split the leaf giving the greatest loss reduction
    \EndWhile
    \State Update $F_b(x)\gets F_{b-1}(x)+\eta T_b(x)$
\EndFor
\State Use 500 rounds and minimum child size 5 for GF
\State Use 400 rounds, 7 leaves, $\eta=0.02$ and minimum child size 10 for FC
\State Predict GF rows directly or recursively update the FC target history
\State \Return $F_B(X_Q)$ or the resulting recursive chain
\end{algorithmic}
\end{algorithm}

\begin{algorithm}[H]
\caption{Unweighted voting-regression forecast ensemble}
\label{alg:app-unweighted-voting-regression}
\small
\begin{algorithmic}[1]
\Require Aligned Random Forest, XGBoost and LightGBM forecast chains
\Ensure Equal-weight ensemble forecast $\hat y^{ens}$
\For{each tower, origin and forecast day $t$}
    \State Verify that all component predictions refer to the same target row
    \State $\hat y_t^{ens}\gets
    (\hat y_t^{RF}+\hat y_t^{XGB}+\hat y_t^{LGBM})/3$
\EndFor
\State Evaluate the ensemble only at QC-valid observed target rows
\State \Return $\hat y^{ens}$
\end{algorithmic}
\end{algorithm}

\section{Neural Models}

The ANN is retained as a published-baseline replication
\cite{rumelhart1986backprop}. SAITS and the bidirectional LSTM use observations
on both sides of historical gaps, whereas the forecasting LSTM, TFT and DLinear
operate causally beyond each forecast origin
\cite{du2023saits,hochreiter1997lstm,schuster1997bidirectional,lim2021tft,zeng2023}.

\begin{algorithm}[H]
\caption{Feed-forward ANN gap reconstruction}
\label{alg:app-ann}
\small
\begin{algorithmic}[1]
\Require Standardised observed pairs $(X_O,y_O)$; masked queries $X_H$;
network layers $1{:}L$
\Ensure Reconstructed target $\hat y_H$
\State Initialise weights $W^{(\ell)}$ and biases $b^{(\ell)}$
\For{each training epoch}
    \State Set $a^{(0)}\gets X_O$
    \For{$\ell\gets1$ to $L-1$}
        \State $a^{(\ell)}\gets\phi(W^{(\ell)}a^{(\ell-1)}+b^{(\ell)})$
    \EndFor
    \State $\hat y\gets W^{(L)}a^{(L-1)}+b^{(L)}$
    \State Back-propagate prediction loss and update all trainable parameters
\EndFor
\State Apply the fitted network to $X_H$ and reverse target scaling
\State \Return $\hat y_H$
\end{algorithmic}
\end{algorithm}

\begin{algorithm}[H]
\caption{SAITS diagonally masked attention reconstruction}
\label{alg:app-saits}
\small
\begin{algorithmic}[1]
\Require Per-tower sequence $Z=[X,y]$ with $y_H$ masked; 336-hour windows;
24-hour stride
\Ensure Reconstructed target $\hat y_H$
\State Remove the eight explicit soil-lag channels and standardise the masked sequence
\State Create overlapping windows and a chronological training--validation split
\State Initialise 3 layers, $d_{model}=256$, 4 heads, $d_k=d_v=64$,
$d_{ff}=512$, dropout 0.1
\For{epoch $e\gets1$ to $100$}
    \State Compute diagonally masked attention
    $A=\operatorname{softmax}(QK^\top/\sqrt{d_k}+D)V$, where $D_{tt}=-\infty$
    \State Obtain reconstructions $\widetilde Z^{(1)}$ and $\widetilde Z^{(2)}$
    \State Combine them as
    $\widetilde Z=\beta\odot\widetilde Z^{(1)}+(1-\beta)\odot\widetilde Z^{(2)}$
    \State Minimise spike-weighted absolute reconstruction error in batches of 32
    \State Retain the best validation state; stop after 10 non-improving epochs
\EndFor
\State Average overlapping estimates, reverse scaling and \Return values at $H$
\end{algorithmic}
\end{algorithm}

\begin{algorithm}[H]
\caption{Bidirectional LSTM gap reconstruction}
\label{alg:app-bilstm}
\small
\begin{algorithmic}[1]
\Require Standardised 336-hour windows $[X,y]$; observation mask $m$;
held-out set $H$
\Ensure Reconstructed target $\hat y_H$
\State Replace missing target inputs by zero and append mask $m_t$ as an input channel
\State Initialise a two-layer Bi-LSTM with 128 hidden units per direction and dropout 0.1
\For{epoch $e\gets1$ to $50$}
    \State Randomly mask 20\% of observed targets to form reconstruction mask $\widetilde m$
    \State $\overrightarrow h_t\gets\operatorname{LSTM}(u_t,\overrightarrow h_{t-1})$
    \State $\overleftarrow h_t\gets\operatorname{LSTM}(u_t,\overleftarrow h_{t+1})$
    \State $\hat y_t\gets W[\overrightarrow h_t;\overleftarrow h_t]+b$
    \State Minimise
    $\sum_t\widetilde m_t(1+|y_t|)|y_t-\hat y_t|/\sum_t\widetilde m_t$
    using Adam and learning rate $10^{-3}$
    \State Retain the best validation state; stop after 8 non-improving epochs
\EndFor
\State Average overlapping estimates, reverse scaling and \Return values at $H$
\end{algorithmic}
\end{algorithm}

\begin{algorithm}[H]
\caption{Causal LSTM forecasting}
\label{alg:app-lstm-forecast}
\small
\begin{algorithmic}[1]
\Require Pre-origin daily target and drivers; 28-day lookback; hidden size 64
\Ensure Continuous 365-day forecast
\State Form training windows ending no later than the forecast origin
\State Standardise using pre-origin training statistics
\For{each training window and time $t$}
    \State Compute input, forget and output gates $i_t,f_t,o_t$
    \State $\widetilde c_t\gets\tanh(W_cu_t+U_ch_{t-1}+b_c)$
    \State $c_t\gets f_t\odot c_{t-1}+i_t\odot\widetilde c_t$
    \State $h_t\gets o_t\odot\tanh(c_t)$ and map $h_t$ to the next target
    \State Update parameters by back-propagation through time
\EndFor
\For{forecast day $h\gets1$ to $365$}
    \State Predict from the latest 28 permitted days and known day-$h$ drivers
    \State Append the prediction, rather than the later observation, to the history
\EndFor
\State Reverse target scaling and \Return the forecast chain
\end{algorithmic}
\end{algorithm}

\begin{algorithm}[H]
\caption{Temporal Fusion Transformer forecasting}
\label{alg:app-tft}
\small
\begin{algorithmic}[1]
\Require Pre-origin daily sequences; known-future covariates; 28-day lookback;
$d_{model}=32$; 4 heads; dropout 0.1
\Ensure Continuous 365-day forecast and optional quantiles
\State Encode static, observed-history and known-future variables separately
\State Apply variable-selection networks to obtain time-varying feature weights
\State Pass selected inputs through gated residual networks and an LSTM encoder--decoder
\State Apply causal temporal attention over permitted positions
\State Produce point or quantile outputs through the gated output head
\For{forecast day $h\gets1$ to $365$}
    \State Supply the latest 28-day context and covariates known for day $h$
    \State Retain the next prediction and append it to the target history
\EndFor
\State \Return the continuous chain and any requested quantiles
\end{algorithmic}
\end{algorithm}

\begin{algorithm}[H]
\caption{DLinear decomposition forecasting}
\label{alg:app-dlinear}
\small
\begin{algorithmic}[1]
\Require Pre-origin daily sequence $z_{t-27:t}$; moving-average kernel 25
\Ensure Continuous 365-day forecast
\State Estimate trend $\tau\gets\operatorname{MovingAverage}_{25}(z)$
\State Calculate residual component $s\gets z-\tau$
\State Learn linear projections $W_\tau,b_\tau$ and $W_s,b_s$ from training windows
\State Define $\hat z\gets(W_\tau\tau+b_\tau)+(W_ss+b_s)$
\For{forecast day $h\gets1$ to $365$}
    \State Predict from the latest 28-day sequence
    \State Append the prediction and advance the input window by one day
\EndFor
\State Reverse target scaling and \Return the continuous chain
\end{algorithmic}
\end{algorithm}

\section{Statistical and Kernel Models}

Support-vector regression is retained as a literature-aligned gap-filling
replication \cite{drucker1997svr}. SARIMAX provides the forecasting benchmark
with an explicit stochastic time-series structure \cite{box2015timeseries}.

\begin{algorithm}[H]
\caption{$\epsilon$-support-vector regression gap reconstruction}
\label{alg:app-svr}
\small
\begin{algorithmic}[1]
\Require Standardised training pairs $(X_O,y_O)$; queries $X_H$; kernel $K$;
penalty $C$; tube width $\epsilon$
\Ensure Reconstructed target $\hat y_H$
\State Solve
$\min_{w,b,\xi,\xi^*}\frac12\lVert w\rVert^2+C\sum_i(\xi_i+\xi_i^*)$
subject to $|y_i-(w^\top\phi(x_i)+b)|\le\epsilon+\xi_i$
\State Retain rows with non-zero dual coefficients as support vectors
\State Define
$\hat y(x)\gets\sum_{i\in O}(\alpha_i-\alpha_i^*)K(x_i,x)+b$
\State Reverse target scaling and \Return $\hat y(X_H)$
\end{algorithmic}
\end{algorithm}

\begin{algorithm}[H]
\caption{AIC-selected SARIMAX forecasting}
\label{alg:app-sarimax}
\small
\begin{algorithmic}[1]
\Require Per-tower pre-origin target $y_t$; exogenous vector $x_t$;
$p\in\{1,2\}$, $d=1$, $q\in\{0,1\}$
\Ensure Continuous 365-day forecast
\For{each candidate order $(p,1,q)$}
    \State Fit
    $\Phi_p(B)(1-B)y_t=c+\beta^\top x_t+\Theta_q(B)\varepsilon_t$
    by maximum likelihood on pre-origin observations
    \State Calculate $\mathrm{AIC}=2k-2\log\widehat L$
\EndFor
\State Select the converged candidate with minimum AIC
\State Forecast 365 steps using its dynamic state update and known future $x_t$
\State Propagate the state without inserting later observed FCH\textsubscript{4}
\State \Return predictive means and state-space intervals
\end{algorithmic}
\end{algorithm}

\section{Pretrained Tabular Models}

\subsection{Pretraining and in-context adaptation}

TabPFN and TabICLv2 are prior-data fitted networks (PFNs).  Their model
parameters are learned once, before application to this study, across a large
distribution of synthetic tabular prediction tasks.  For a synthetic task
$\mathcal T\sim p_{\mathrm{synth}}$, rows are divided into a labelled context
$\mathcal D_C=(X_C,y_C)$ and an unlabelled query $X_Q$ with held-out target
$y_Q$.  Pretraining minimises the query negative log-likelihood,
\begin{equation}
\theta^*=\arg\min_\theta
\mathbb E_{\mathcal T\sim p_{\mathrm{synth}}}
\left[-\sum_{j\in Q}\log q_\theta
\left(y_j\mid x_j,\mathcal D_C\right)\right],
\label{eq:app-pfn-pretraining}
\end{equation}
so that one forward pass subsequently approximates the posterior predictive
rule induced by the synthetic task prior
\cite{hollmann2023tabpfn,hollmann2025tabpfn,qu2025tabicl,qu2026tabiclv2}.
This is amortised inference: the expensive learning occurs across synthetic
tasks during offline pretraining, rather than through gradient optimisation on
each new dataset.

TabPFN v2 retains a representation for every table cell.  Numerical values are
combined with feature-specific random embeddings, grouped in pairs, and passed
through 12 Transformer blocks (192-dimensional embeddings and six attention
heads).  Each block alternates attention across features within a row and
across samples within a feature.  Context targets are embedded while query
targets remain masked, allowing queries to attend to labelled examples without
target leakage.  For regression, the output head assigns probability to target
intervals through a full-support bar distribution and is trained by negative
log-likelihood; means, medians or quantiles are then derived from this
distribution \cite{hollmann2025tabpfn}.

Its synthetic prior samples table properties and structural causal models:
random directed acyclic graphs are populated with nonlinear mechanisms,
transformations, feature types, missingness and noise.  Approximately 100
million such tasks expose the model to varied tabular relationships.  This
study used the released generic \texttt{tabpfn-v2-\allowbreak regressor.ckpt}, not a
time-series-fine-tuned checkpoint.

TabICLv2 instead compresses features before dataset-level inference.  Repeated
circular feature groups are given target-aware embeddings for labelled rows; a
column Transformer processes each feature across samples, and a row Transformer
combines the features into a 512-dimensional representation.  A 12-layer
dataset Transformer then relates labelled and query rows.  Its query-aware
scalable softmax (QASSMax) counters attention dilution as the context grows,
reducing the principal scaling from approximately
$O(n^2m+nm^2)$ for cell-level TabPFN to $O(n^2+nm^2)$
\cite{qu2025tabicl,qu2026tabiclv2}.

TabICLv2 is pretrained on approximately 35 million filtered synthetic tasks
generated by random graphs containing neural, tree, Gaussian-process, linear,
quadratic, product and discretisation mechanisms.  Its three-stage curriculum
increases contexts from 1,024 to at most 60,000 rows and 100 features.  The
regression model uses Muon optimisation and pinball loss over 999 conditional
quantiles.  Thus, unlike squared-error regressors, neither TabPFN nor TabICLv2
directly optimises the downstream $R^2$, RMSE, MAE or MASE metrics.

Neither checkpoint was pretrained on North Wyke Farm Platform data, methane
fluxes or the held-out forecast years.  Here, \texttt{fit()} only preprocesses
and registers the chronologically permitted labelled context; it does not
update $\theta^*$.  Temporal information therefore enters through context
selection and the engineered calendar, environmental and management features.
The direct forecasting configuration used eight internal estimators (TabPFN
seed 137; TabICLv2 seed 42) without checkpoint fine-tuning.  Forecast
combinations follow Bates and Granger's variance-reduction principle
\cite{bates1969combination}.

\begin{algorithm}[H]
\caption{Published offline pretraining scheme for TabPFN and TabICLv2}
\label{alg:app-pfn-pretraining}
\small
\begin{algorithmic}[1]
\Require Model family $M\in\{\mathrm{TabPFN},\mathrm{TabICLv2}\}$;
synthetic task prior $p_{\mathrm{synth}}^M$; model $q_\theta$
\Ensure Released pretrained checkpoint $\theta^*$
\For{each offline pretraining step}
    \State Sample table dimensions, feature types and task difficulty
    \If{$M=\mathrm{TabPFN}$}
        \State Sample an SCM directed acyclic graph, node mechanisms and noise
    \Else
        \State Sample a random graph with mixed neural, tree, GP and algebraic mechanisms
    \EndIf
    \State Generate a synthetic table $\mathcal T=(X,y)$ and apply model-specific filtering
    \State Split rows into labelled context $\mathcal D_C$ and query $(X_Q,y_Q)$
    \State Predict $q_\theta(y_Q\mid X_Q,\mathcal D_C)$ in one forward pass
    \If{$M=\mathrm{TabPFN}$}
        \State Compute the discretised regression negative log-likelihood and update with Adam
    \Else
        \State Sum pinball losses over 999 quantiles and update with Muon
        \State Advance the three-stage table-size curriculum
    \EndIf
\EndFor
\State Freeze the final parameters; no study data enter this procedure
\State \Return $\theta^*$ for subsequent in-context inference
\end{algorithmic}
\end{algorithm}

\begin{algorithm}[H]
\caption{TabPFN inference for gap filling or direct daily forecasting}
\label{alg:app-tabpfn}
\small
\begin{algorithmic}[1]
\Require Labelled context $\mathcal D_C=\{(x_i,y_i)\}$; query features $X_Q$;
pretrained parameters $\theta$; task $\tau$
\Ensure Predictive distribution and point predictions
\If{$\tau=\mathrm{GF}$}
    \State Construct an hourly tabular task after masking $y_H$ and exclude $H$ from context
\Else
    \State Use pre-origin daily rows and the 52-driver plus 5 tower/time-indicator frame
    \State Optionally restrict context to the most recent 1,095 or 1,460 days
\EndIf
\State Fit training-only missing-value preprocessing and configured target scaling
\State Jointly encode labelled context and unlabelled query rows with the pretrained model
\State Infer $p_\theta(y_Q\mid\mathcal D_C,X_Q)$ without gradient-based task training
\State Extract the configured point statistic (median for direct forecasts) and quantiles
\State Reverse target scaling and \Return predictions
\end{algorithmic}
\end{algorithm}

\begin{algorithm}[H]
\caption{TabICLv2 inference for gap filling or direct daily forecasting}
\label{alg:app-tabicl}
\small
\begin{algorithmic}[1]
\Require Labelled context $\mathcal D_C$; query features $X_Q$;
pretrained parameters $\theta$; task $\tau$
\Ensure Point predictions and optional quantiles
\If{$\tau=\mathrm{GF}$}
    \State Train separately by tower on the expanded 30-predictor frame
    \State Exclude $H$ and sample at most 10,000 observed rows using seed 42
    \State Mean-impute predictors using the sampled context
\Else
    \State Use pre-origin daily context and the applicable forecasting frame
    \State Optionally retain all history, 1,095 days or 1,460 days
    \State Apply tower-wise target scaling for the configured 1,460-day variant
\EndIf
\State Register the labelled context through the model fit interface
\State Keep the pretrained checkpoint fixed; \texttt{fit()} performs no gradient update
\State Infer $p_\theta(y_Q\mid\mathcal D_C,X_Q)$ jointly from context and queries
\State Extract the median point prediction and optional 0.05, 0.50 and 0.95 quantiles
\State Reverse target scaling and \Return predictions
\end{algorithmic}
\end{algorithm}

\begin{algorithm}[H]
\caption{Probability-gated spike correction for a foundation-model forecast}
\label{alg:app-spike-correction}
\small
\begin{algorithmic}[1]
\Require Base regressor $f_b$; spike classifier $c$; spike regressor $f_s$;
normal regressor $f_n$; pre-origin tower-specific 95th percentile $q_{0.95}$
\Ensure Corrected forecast $\hat y^{corr}$
\State Label context row $i$ as a spike if $y_i>q_{0.95}$
\State Fit or condition $c$, $f_s$, $f_n$ and $f_b$ using permitted context only
\For{each query day $t$}
    \State $p_t^{spike}\gets c(x_t)$
    \State $e_t\gets[f_s(x_t)-f_n(x_t)]_+$
    \State $\hat y_t^{corr}\gets f_b(x_t)+0.25\,p_t^{spike}e_t$
\EndFor
\State \Return $\hat y^{corr}$
\end{algorithmic}
\end{algorithm}

\begin{algorithm}[H]
\caption{Best-performing TabICL forecasting model}
\label{alg:app-best-tabicl-forecast}
\small
\begin{algorithmic}[1]
\Require TabICLv2 forecaster $F$; pre-origin daily data; 365-day queries
\Ensure Final blended forecast $\hat y^{\mathrm{best}}$
\State $\hat y^{(1)}\gets$ all-history forecast from $F$ with
Algorithm~\ref{alg:app-spike-correction}
\State $\hat y^{(2)}\gets$ forecast from $F$ using the latest 1,095 context days
\State $\hat y^{(3)}\gets$ forecast from $F$ using the latest 1,460 context days and
tower-wise target scaling
\For{each forecast day $t$}
    \State $\hat y_t^{\mathrm{best}}\gets
    (\hat y_t^{(1)}+\hat y_t^{(2)}+\hat y_t^{(3)})/3$
\EndFor
\State Retain the three component chains and evaluate the blend on identical observed rows
\State \Return $\hat y^{\mathrm{best}}$
\end{algorithmic}
\end{algorithm}

\FloatBarrier

